% !TEX root = chapter-standalone.tex

\section{Twisted arrow construction}
\todotext{\alphubel: @JL/@Gioele: Don't use \str{stylemorph}, use the various $\morc,\mord,\more,\etc.$
    It's very hard to edit, otherwise.
    Don't forget to change the graphics construction.
}

\todotext{@J: review and edit}

The arrow construction of the previous section admits two useful variants, in which one of the two connecting morphisms is reversed---``twisted''.
Objects are still the morphisms of~\CatC.
In the \emph{source-twisted} variant, a morphism from~$\mora$ to~$\morb$ consists of a morphism between sources going \emph{backwards} and a morphism between targets going forwards; in the \emph{target-twisted} variant, it is the other way around.

One intuition for the source-twisted variant: a morphism from~$\mora$ to~$\morb$ witnesses that~$\morb$ ``extends''~$\mora$, in the sense that~$\mora$ occurs as the middle portion of~$\morb$, with some pre-processing prepended and some post-processing appended.
In the target-twisted variant, the roles are swapped: a morphism from~$\mora$ to~$\morb$ witnesses that~$\mora$ extends~$\morb$.
These constructions appear naturally when studying tradeoffs and interval-like structures, as the example below illustrates.

\begin{ctdefinition}[Source-twisted arrow category]
    \label{def:twisted-arrow-category}
    \label{def:twisted-arrow}
    \label{def:source-twisted-arrow}
    \SYNDEF{twisted arrow category}
    Given a category \CatC, its \maindef{source-twisted arrow category} $\TwistedArrow\CatC$ is the category given by:

    \constit

    \begin{enumerate}
        \item Objects: morphisms in \CatC.
        \item Morphisms:
              A morphism in $\TwistedArrow\CatC$ from $\mora\colon \Obja \mtoin{\CatC} \Objb $ to $\morb \colon \Objc \mtoin{\CatC} \Objd$ is given by a pair of morphisms $\twarrowmorasource \colon \Objc \mto \Obja$ and $\twarrowmoratarget \colon \Objb \mto \Objd$ in $\CatC$ such that the following diagram commutes:
              \begin{equation}\label{eq:180_twistedarrow}
                  \includesag{180_twistedarrow}
              \end{equation}
              that is, such that $\twarrowmorasource \mthen \mora \mthen \twarrowmoratarget = \morb$.
        \item Composition:
              Composition in~$\TwistedArrow\CatC$ is given by placing \SY{commutative squares} side by side.
              Consider~$\tup{\twarrowmorasource,\twarrowmoratarget} \colon \mora \mto \morb$ and~$\tup{\twarrowmorbsource,\twarrowmorbtarget} \colon \morb \mto \morc$ in~$\TwistedArrow\CatC$, giving rise to the following composite diagram
              \begin{equation}\label{eq:180_twistedarrow_comp}
                  \includesag{180_twistedarrow_comp}.
              \end{equation}
              It is commutative because its component squares are, and hence it defines a morphism $\mora \mto \morc$ in $\TwistedArrow\CatC$, namely
              \begin{equation}\label{eq:twarrow-cat-composition}
                  \tup{\twarrowmorasource,\twarrowmoratarget}\mthenof{\TwistedArrow \CatC} \tup{\twarrowmorbsource,\twarrowmorbtarget}
                  \definedas \tup{ \twarrowmorbsource \mthenof\CatC \twarrowmorasource, \twarrowmoratarget\mthenof\CatC \twarrowmorbtarget }.
              \end{equation}
        \item Identity morphisms: given an object $\mora \colon \Obja \mto \Objb$ of $\TwistedArrow\CatC$, its identity morphism is $\tup{\catid_{\Obja},\catid_{\Objb}}$.
    \end{enumerate}

    When we need to distinguish it from the target-twisted variant below, we also write~$\TwistedArrowSource\CatC$ for~$\TwistedArrow\CatC$.
\end{ctdefinition}

\begin{ctdefinition}[Target-twisted arrow category]
    \label{def:target-twisted-arrow}
    Given a category \CatC, its \maindef{target-twisted arrow category} $\TwistedArrowTarget\CatC$ is the category given by:

    \constit

    \begin{enumerate}
        \item Objects: morphisms in \CatC.
        \item Morphisms:
              A morphism in $\TwistedArrowTarget\CatC$ from $\mora\colon \Obja \mtoin{\CatC} \Objb $ to $\morb \colon \Objc \mtoin{\CatC} \Objd$ is given by a pair of morphisms $\twarrowmorasource \colon \Obja \mto \Objc$ and $\twarrowmoratarget \colon \Objd \mto \Objb$ in $\CatC$ such that the following diagram commutes:
              \begin{equation}\label{eq:target-twistedarrow}
                  \begin{tikzcd}[every arrow/.append style={arrowmorphismstyle}]
                      \Obja \arrow[r, "\mora"] \arrow[d, "\twarrowmorasource"'] \pgfmatrixnextcell \Objb \\
                      \Objc \arrow[r, "\morb"'] \pgfmatrixnextcell \Objd \arrow[u, "\twarrowmoratarget"']
                  \end{tikzcd}
              \end{equation}
              that is, such that $\twarrowmorasource \mthen \morb \mthen \twarrowmoratarget = \mora$.
        \item Composition:
              Given~$\tup{\twarrowmorasource,\twarrowmoratarget} \colon \mora \mto \morb$ and~$\tup{\twarrowmorbsource,\twarrowmorbtarget} \colon \morb \mto \morc$ in~$\TwistedArrowTarget\CatC$, we define
              \begin{equation}\label{eq:target-twarrow-cat-composition}
                  \tup{\twarrowmorasource,\twarrowmoratarget}\mthenof{\TwistedArrowTarget \CatC} \tup{\twarrowmorbsource,\twarrowmorbtarget}
                  \definedas \tup{ \twarrowmorasource \mthenof\CatC \twarrowmorbsource, \twarrowmorbtarget\mthenof\CatC \twarrowmoratarget }.
              \end{equation}
              This is again a morphism of~$\TwistedArrowTarget\CatC$, since~$\twarrowmorasource \mthen \twarrowmorbsource \mthen \morc \mthen \twarrowmorbtarget \mthen \twarrowmoratarget = \twarrowmorasource \mthen \morb \mthen \twarrowmoratarget = \mora$.
        \item Identity morphisms: given an object $\mora \colon \Obja \mto \Objb$ of $\TwistedArrowTarget\CatC$, its identity morphism is $\tup{\catid_{\Obja},\catid_{\Objb}}$.
    \end{enumerate}
\end{ctdefinition}

\begin{remark}[How the two twisted constructions are related]
    \label{rem:source-target-twisted}
    Compare the two definitions.
    A morphism from~$\mora$ to~$\morb$ in~$\TwistedArrowTarget\CatC$ is a pair~$\twarrowmorasource \colon \Obja \mto \Objc$, $\twarrowmoratarget \colon \Objd \mto \Objb$ with~$\twarrowmorasource \mthen \morb \mthen \twarrowmoratarget = \mora$.
    But this is \emph{exactly} the data, subject to exactly the same condition, of a morphism from~$\morb$ to~$\mora$ in~$\TwistedArrowSource\CatC$.
    Composition and identities also match up, so the target-twisted arrow category is the opposite of the source-twisted one:
    \begin{equation}
        \TwistedArrowTarget\CatC \simeq \pars{\TwistedArrowSource\CatC}\op.
    \end{equation}
    In this sense, there is really only one new construction here, together with its opposite.
    This is why, in the literature, one usually finds only ``the'' twisted arrow category---but beware that different authors take different variants as their default.
    In this book, $\TwistedArrow\CatC$ without further qualification means the source-twisted variant~$\TwistedArrowSource\CatC$.

    Together with \cref{rem:arrow-opposite}, this gives a complete picture of the four ways of choosing directions for the two connecting morphisms between~$\mora$ and~$\morb$:
    \begin{center}
        \begin{tabular}{lll}
            source leg & target leg & category \\
            \hline
            forwards   & forwards   & $\Arrow\CatC$ \\
            backwards  & backwards  & $\pars{\Arrow\CatC}\op \simeq \Arrow\pars{\CatCop}$ \\
            backwards  & forwards   & $\TwistedArrowSource\CatC$ \\
            forwards   & backwards  & $\TwistedArrowTarget\CatC \simeq \pars{\TwistedArrowSource\CatC}\op$
        \end{tabular}
    \end{center}
\end{remark}

\begin{gradedexercise}[\exname{TwistedCat}]
    \label{ex:TwistedCat}

    Let $\CatC$ be a category, and let $\TwistedArrow\CatC$ be the associated twisted arrow category.
    Check that the definition of $\TwistedArrow\CatC$ does indeed define a category.
    Specifically, check that for $\TwistedArrow\CatC$
    \begin{enumerate}
        \item composition of composable morphism does again define a morphism of $\TwistedArrow\CatC$;
        \item composition is associative;
        \item identity morphisms satisfy the identity laws (that they behave neutrally for composition).
    \end{enumerate}
\end{gradedexercise}

\solutionof{TwistedCat}

\begin{example}[Twisted construction in posets and categories]
    \label{exa:twisted-arrow-poset}
    Consider a poset~\posA.
    The twisted arrow category~$\TwistedArrow \poscat{\posA}$ is isomorphic to the \SY{poset} (viewed as a category)~$\poscat{\pars{\posint\posA}}$ of nonempty \emph{intervals} in~\posA:
    \begin{equation}
        \TwistedArrow (\poscat \posA) \simeq \poscat{(\posint\posA)}.
    \end{equation}
    \todotext{\alphubel: Did we define what is an isomorphic category, at this point?}
\end{example}

\begin{exercise}
    \label{ex:twisted-arrow-poset}
    Prove the statement in \cref{exa:twisted-arrow-poset}.
    Recall (\cref{def:interval}) that an interval~$\interv{\posela}{\poselb}$ of~\posA is a pair of elements with~$\posela \posAleq \poselb$, with associated subset
    \begin{equation}
        \interv{\posela}{\poselb}
        \definedas \makeset{\poselc \setin \posAset \mid \posela \posAleq \poselc \posAleq \poselb}.
    \end{equation}
    Start by showing that the objects of~$\TwistedArrow (\poscat \posA)$ correspond to the intervals of~\posA.
    Then show that there is a morphism between two such objects exactly when the corresponding intervals are related in~$\posint\posA$ (\cref{def:twisted-poset-of-intervals}).
\end{exercise}
\begin{solution}
    \emph{Objects.}
    An object of~$\TwistedArrow (\poscat \posA)$ is a morphism of~$\poscat{\posA}$, that is, a pair of elements~$\posela, \poselb \setin \posAset$ with~$\posela \posAleq \poselb$ (between two elements there is at most one morphism).
    This is exactly an interval~$\interv{\posela}{\poselb}$ of~\posA, that is, an element of~$\posint\posA$.

    \emph{Morphisms.}
    Consider two objects~$\mora \colon \posela \mto \poselb$ and~$\morb \colon \poselc \mto \poseld$.
    A morphism~$\mora \mto \morb$ in~$\TwistedArrow (\poscat \posA)$ is a pair consisting of a morphism~$\poselc \mto \posela$ and a morphism~$\poselb \mto \poseld$ of~$\poscat{\posA}$ such that the defining square commutes.
    Such a pair exists if and only if~$\poselc \posAleq \posela$ and~$\poselb \posAleq \poseld$.
    Moreover, since~$\poscat{\posA}$ has at most one morphism between any two objects, the square commutes automatically, and the pair is unique when it exists.
    So there is exactly one morphism~$\mora \mto \morb$ if~$\poselc \posAleq \posela$ and~$\poselb \posAleq \poseld$, and none otherwise.
    By \cref{eq:twisted-2}, this is exactly the condition~$\interv{\posela}{\poselb} \posleqof{\posint\posA} \interv{\poselc}{\poseld}$.

    \emph{Composition and identities.}
    Both categories have at most one morphism between any two objects, so composition and \SY{identity morphisms} necessarily agree.
    Hence the two categories are the same up to renaming objects:~$\TwistedArrow (\poscat \posA) \simeq \poscat{(\posint\posA)}$.
\end{solution}

\devel{
    \begin{remark}
        Recall \cref{sec:posets-as-cats} and note that the map which sends a \SY{poset} (a category) to its twisted arrow category is a \SY{functor}, which sends objects of the poset
        \todotextjira{116}{\alphubel: @JL: Finish remark about twisted map functor}
    \end{remark}
}
